\chapter{Matériaux et sections}

\section{Matériaux}

Le logiciel contient une bibliothèque de matériaux usuels :

\begin{itemize}
  \item béton selon les classes Eurocode 2 ;
  \item acier selon les nuances Eurocode 3 ;
  \item armatures pour les vérifications futures.
\end{itemize}

Chaque matériau possède des propriétés mécaniques telles que le module d'Young,
le coefficient de Poisson, la masse volumique ou les résistances de calcul.

\section{Sections de barre}

Les sections de barre décrivent les caractéristiques géométriques des éléments
linéaires :

\begin{itemize}
  \item aire ;
  \item inerties ;
  \item torsion ;
  \item matériau associé.
\end{itemize}

Les sections disponibles incluent notamment :

\begin{itemize}
  \item sections rectangulaires ;
  \item sections en T ;
  \item profilés acier européens IPE, HEA et HEB.
\end{itemize}

\section{Sections plaque}

Les sections plaque définissent principalement :

\begin{itemize}
  \item l'épaisseur ;
  \item le matériau ;
  \item la formulation de l'élément surfacique.
\end{itemize}

\section{Bonnes pratiques}

\begin{itemize}
  \item Nommer les sections clairement, par exemple \fichier{Poteau 30x30}
  ou \fichier{Dalle 20 cm}.
  \item Vérifier les unités avant d'ajouter une section.
  \item Utiliser une bibliothèque de sections stable pour éviter les incohérences
  entre plusieurs modèles.
\end{itemize}
